\documentclass[11pt,landscape,letterpaper]{article}

\usepackage[margin=0.9cm]{geometry}
\usepackage[utf8]{inputenc}
\usepackage[T1]{fontenc}
\usepackage[spanish]{babel}
\usepackage{amsmath,amssymb}
\usepackage{adjustbox}
\usepackage{tikz}
\usetikzlibrary{arrows.meta,positioning,shapes.geometric,calc}

\pagestyle{empty}

\definecolor{azul}{RGB}{31,78,121}
\definecolor{azulclaro}{RGB}{221,235,247}
\definecolor{verde}{RGB}{226,239,218}
\definecolor{amarillo}{RGB}{255,242,204}
\definecolor{gris}{RGB}{242,242,242}

\tikzset{
    head/.style={
        rectangle,
        rounded corners=8pt,
        draw=azul,
        fill=azulclaro,
        very thick,
        text width=7.0cm,
        minimum height=1.2cm,
        align=center,
        font=\bfseries\footnotesize
    },
    procbox/.style={
        rectangle,
        rounded corners=4pt,
        draw=black!70,
        fill=gris,
        thick,
        text width=5.8cm,
        minimum height=1.1cm,
        align=center,
        font=\footnotesize
    },
    steplist/.style={
        procbox,
        text width=8.3cm,
        align=left
    },
    formula/.style={
        rectangle,
        rounded corners=4pt,
        draw=azul!70!black,
        fill=azulclaro!50,
        thick,
        text width=6.9cm,
        minimum height=1.1cm,
        inner sep=4pt,
        align=center,
        font=\scriptsize
    },
    output/.style={
        rectangle,
        rounded corners=4pt,
        draw=black!70,
        fill=verde,
        thick,
        text width=6.3cm,
        minimum height=1.1cm,
        align=center,
        font=\footnotesize
    },
    decision/.style={
        diamond,
        aspect=2.25,
        draw=black!70,
        fill=amarillo,
        thick,
        text width=4.8cm,
        inner sep=2pt,
        align=center,
        font=\bfseries\footnotesize
    },
    note/.style={
        rectangle,
        rounded corners=4pt,
        draw=azul,
        fill=azul!8,
        thick,
        text width=21.0cm,
        minimum height=1.1cm,
        align=left,
        font=\scriptsize
    },
    arrow/.style={
        -{Latex[length=2.5mm]},
        thick,
        draw=black!75,
        rounded corners=4pt
    },
    sidearrow/.style={
        -{Latex[length=2.2mm]},
        thick,
        draw=azul!80!black,
        rounded corners=4pt
    }
}

\begin{document}

\noindent\makebox[\textwidth][c]{%
\begin{adjustbox}{max totalsize={\textwidth}{0.98\textheight},center}
\begin{tikzpicture}[x=0.78cm,y=0.17cm]

% =========================================================
% BLOQUE SUPERIOR
% =========================================================

\node[head] (inicio) at (0,0) {
Inicio\\
Obtener la informaci\'on financiera de cuatro empresas p\'ublicas
};

\node[procbox] (empresas) at (0,-8) {
Empresas elegidas\\
dos del mercado Asi\'atico y dos de mercados europeos
};

\node[steplist] (datos) at (0,-18) {
\textbf{Integrar insumos financieros}\\
capital de trabajo, utilidades retenidas, EBIT, valor de mercado del patrimonio neto, ventas, activo total, pasivo total, pasivos de largo plazo
};

\node[decision] (ramas) at (0,-31) {
Determinar insumos de riesgo
};

% =========================================================
% RAMA ALTMAN
% =========================================================

\node[procbox] (altman0) at (-20,-46) {
Z de Altman\\
calidad crediticia y probabilidad de incumplimiento
};

\node[procbox] (altman1t) at (-20,-59) {
Obtener razones financieras
};
\node[formula] (altman1m) at (-8,-59) {
\(
\begin{aligned}
X_1 &= \frac{\text{capital de trabajo}}{\text{activo total}}\\
X_2 &= \frac{\text{utilidades retenidas}}{\text{activo total}}\\
X_3 &= \frac{\text{EBIT}}{\text{activo total}}\\
X_4 &= \frac{\text{valor de mercado del patrimonio}}{\text{pasivo total}}\\
X_5 &= \frac{\text{ventas}}{\text{activo total}}
\end{aligned}
\)
};

\node[procbox] (altman2t) at (-20,-73) {
Calcular puntuaci\'on discriminante
};
\node[formula] (altman2m) at (-8,-73) {
\(
\begin{aligned}
Z &= 1.2X_1+1.4X_2+3.3X_3+0.60X_4+1.00X_5\\
Z' &= 0.717X_1+0.847X_2+3.107X_3+0.420X_4'+0.998X_5\\
Z'' &= 3.25+6.56X_1+3.26X_2+6.72X_3+1.05X_4'
\end{aligned}
\)
};

\node[procbox] (altman3t) at (-20,-87) {
Clasificar zonas discriminantes
};
\node[formula] (altman3m) at (-8,-87) {
\(
\begin{aligned}
Z &> 2.99 &&\Rightarrow \text{Zona segura}\\
1.81 &< Z < 2.99 &&\Rightarrow \text{Zona gris}\\
Z &< 1.81 &&\Rightarrow \text{Zona de incumplimiento}
\end{aligned}
\)
};

\node[procbox] (altman4t) at (-20,-101) {
Mercados emergentes\\
usar \(Z''\), m\'etodo BRE y transferir la calificaci\'on crediticia a una \(PD\)
};
\node[formula] (altman4m) at (-8,-101) {
\begin{minipage}{6.1cm}\centering
\scriptsize
Salida intermedia:\\
\(Z,\ Z',\ Z''\)\\
rating sint\'etico\\
\(PD_{\text{Altman}}\)
\end{minipage}
};

\node[output] (altman5) at (-20,-115) {
Salida Altman\\
\(Z\), rating sint\'etico, calidad crediticia, \(PD_{\text{Altman}}\)
};

% =========================================================
% RAMA MERTON
% =========================================================

\node[procbox] (merton0) at (8,-46) {
Modelo estructural de Merton\\
probabilidad de incumplimiento
};

\node[procbox] (merton1t) at (8,-59) {
Obtener insumos
};
\node[formula] (merton1m) at (21,-59) {
\begin{minipage}{6.1cm}\raggedright
\scriptsize
1. Valor de los pasivos de largo plazo.\\
2. Valor de mercado del capital de la empresa.\\
3. Tasa libre de riesgo.\\
4. \'Ultimos 500 precios de la acci\'on de la empresa.
\end{minipage}
};

\node[procbox] (merton2t) at (8,-73) {
Obtener la volatilidad anual de los rendimientos
};
\node[formula] (merton2m) at (21,-73) {
\(
\begin{aligned}
\sigma_E &= \text{volatilidad anual de los rendimientos}\\
d_1 &= \frac{\ln(V_0/D)+\left(r+\frac{\sigma_V^2}{2}\right)T}{\sigma_V\sqrt{T}}\\
d_2 &= d_1-\sigma_V\sqrt{T}
\end{aligned}
\)
};

\node[procbox] (merton3t) at (8,-87) {
Usar Solver para encontrar \(V_0\) y \(\sigma_V\)
};
\node[formula] (merton3m) at (21,-87) {
\(
\begin{aligned}
E_0 &= V_0N(d_1)-De^{-rT}N(d_2)\\
\sigma_EE_0 &= N(d_1)\sigma_VV_0\\
&\text{minimizar }F(x,y)^2+G(x,y)^2
\end{aligned}
\)
};

\node[procbox] (merton4t) at (8,-101) {
Caracterizar la deuda\\
obtener \(PD\)
};
\node[formula] (merton4m) at (21,-101) {
\(
\begin{aligned}
PD &= N(-d_2)\\
PD_{\text{final}} &= \max\{N(-d_2),0.03\%\}
\end{aligned}
\)
};

\node[procbox] (merton5t) at (8,-115) {
Caracterizar la deuda\\
deuda de mercado y deuda sin riesgo
};
\node[formula] (merton5m) at (21,-115) {
\(
\begin{aligned}
D_{0,M} &= V_0-E_0\\
D_{0,MNR} &= D_Te^{-rT}
\end{aligned}
\)
};

\node[procbox] (merton6t) at (8,-129) {
Caracterizar la deuda\\
p\'erdida esperada
};
\node[formula] (merton6m) at (21,-129) {
\(
\begin{aligned}
EL_{\$} &= D_{0,MNR}-D_{0,M}\\
EL_{\%} &= \frac{D_{0,MNR}-D_{0,M}}{D_{0,MNR}}
\end{aligned}
\)
};

\node[procbox] (merton7t) at (8,-143) {
Caracterizar la deuda\\
\(RR\) y \(LGD\)
};
\node[formula] (merton7m) at (21,-143) {
\(
\begin{aligned}
RR &= 1-\frac{EL}{PD}\\
LGD &= 1-RR
\end{aligned}
\)
};

\node[output] (merton8) at (8,-157) {
Salida Merton\\
\(PD_{\text{Merton}}\), valor de mercado de la deuda, valor presente de la deuda, p\'erdida esperada, \(RR\), \(LGD\)
};

% =========================================================
% INTEGRACION DE CARTERA
% =========================================================

\node[note] (nota) at (0,-171) {
\textbf{Cartera de cr\'editos:} los pasivos de largo plazo de las cuatro empresas. La moneda reportada son USD.
};

\node[procbox] (selectt) at (-6,-184) {
Seleccionar insumos para la cartera
};
\node[formula] (selectm) at (8,-184) {
\begin{minipage}{6.1cm}\centering
\scriptsize
usar \(PD\) y \(LGD\)\\
mayores a cero\\
para todas las empresas
\end{minipage}
};

% =========================================================
% MEDIDAS DE RIESGO
% =========================================================

\node[procbox] (creditot) at (-18,-200) {
Medida de riesgo por cr\'edito
};
\node[formula] (creditom) at (-18,-214) {
\(
\begin{aligned}
EL_i &= Exp_i\times LGD_i\times PD_i\\
UL_i &= Exp_i\times LGD_i\times \sqrt{PD_i(1-PD_i)}\\
\varrho(X_i) &= (1+\theta)\mu_{X_i}\\
\varrho(X_i) &= \mu_{X_i}+\alpha\sigma_{X_i}
\end{aligned}
\)
};

\node[procbox] (carterat) at (0,-200) {
Medida de riesgo para la cartera
};
\node[formula] (carteram) at (0,-214) {
\(
\begin{aligned}
EL_P &= \sum_{i=1}^{N}EL_i\\
UL_P &= \sqrt{\sum_{i=1}^{N}UL_i^2+2\sum_{i<j}^{N}\rho_{ij}UL_iUL_j}
\end{aligned}
\)
};

% =========================================================
% ESCENARIOS DE CORRELACION
% =========================================================

\node[decision] (corr) at (18,-200) {
Escenario de correlaci\'on
};

\node[procbox] (rho0t) at (11,-217) {
Correlaci\'on cero
};
\node[formula] (rho0m) at (11,-231) {
\(
\begin{aligned}
\rho_{ij} &= 0\\
UL_{P,0} &= \sqrt{\sum_{i=1}^{N}UL_i^2}
\end{aligned}
\)
};

\node[procbox] (rho1t) at (25,-217) {
Correlaci\'on con precio de las acciones
};
\node[formula] (rho1m) at (25,-231) {
\begin{minipage}{6.1cm}\centering
\scriptsize
\(\rho_{ij}=\mathrm{corr}(r_i,r_j)\)\\
usar rendimientos accionarios como proxy
\end{minipage}
};

% =========================================================
% SALIDA FINAL
% =========================================================

\node[procbox] (reservast) at (18,-247) {
Reservas y capital econ\'omico
};
\node[formula] (reservasm) at (18,-261) {
\(
\begin{aligned}
\text{Reservas} &= EL_P\\
CaR &= EL_P+\alpha UL_P\\
\text{Capital econ\'omico} &= CaR-EL_P
\end{aligned}
\)
};

\node[head] (fin) at (0,-278) {
Fin\\
Reporte de resultados y diagrama de procedimientos operativos
};

% =========================================================
% COORDENADAS AUXILIARES
% =========================================================

\coordinate (branchL) at (-20,-38);
\coordinate (branchR) at (8,-38);
\coordinate (altLane) at (-29,-171);
\coordinate (merLane) at (29,-171);
\coordinate (noteLeftIn) at (-16,-171);
\coordinate (noteRightIn) at (16,-171);
\coordinate (creditoJoin) at (-9,-200);
\coordinate (carteraJoin) at (9,-200);
\coordinate (corrSplitL) at (15,-208);
\coordinate (corrSplitR) at (21,-208);
\coordinate (resLeftJoin) at (13,-247);
\coordinate (resRightJoin) at (23,-247);

% =========================================================
% FLECHAS PRINCIPALES
% =========================================================

\draw[arrow] (inicio.south) -- (empresas.north);
\draw[arrow] (empresas.south) -- (datos.north);
\draw[arrow] (datos.south) -- (ramas.north);

\draw[arrow] (ramas.south) |- (branchL) -- (altman0.north);
\draw[arrow] (ramas.south) |- (branchR) -- (merton0.north);

\draw[arrow] (altman0.south) -- (altman1t.north);
\draw[arrow] (altman1t.south) -- (altman2t.north);
\draw[arrow] (altman2t.south) -- (altman3t.north);
\draw[arrow] (altman3t.south) -- (altman4t.north);
\draw[arrow] (altman4t.south) -- (altman5.north);

\draw[arrow] (merton0.south) -- (merton1t.north);
\draw[arrow] (merton1t.south) -- (merton2t.north);
\draw[arrow] (merton2t.south) -- (merton3t.north);
\draw[arrow] (merton3t.south) -- (merton4t.north);
\draw[arrow] (merton4t.south) -- (merton5t.north);
\draw[arrow] (merton5t.south) -- (merton6t.north);
\draw[arrow] (merton6t.south) -- (merton7t.north);
\draw[arrow] (merton7t.south) -- (merton8.north);

\draw[sidearrow] (altman1t.east) -- (altman1m.west);
\draw[sidearrow] (altman2t.east) -- (altman2m.west);
\draw[sidearrow] (altman3t.east) -- (altman3m.west);
\draw[sidearrow] (altman4t.east) -- (altman4m.west);

\draw[sidearrow] (merton1t.east) -- (merton1m.west);
\draw[sidearrow] (merton2t.east) -- (merton2m.west);
\draw[sidearrow] (merton3t.east) -- (merton3m.west);
\draw[sidearrow] (merton4t.east) -- (merton4m.west);
\draw[sidearrow] (merton5t.east) -- (merton5m.west);
\draw[sidearrow] (merton6t.east) -- (merton6m.west);
\draw[sidearrow] (merton7t.east) -- (merton7m.west);

\draw[arrow] (altman5.south) |- (altLane) -- (noteLeftIn) -- (nota.west);
\draw[arrow] (merton8.south) |- (merLane) -- (noteRightIn) -- (nota.east);
\draw[arrow] (nota.south) -- (selectt.north);
\draw[sidearrow] (selectt.east) -- (selectm.west);

\draw[arrow] (selectt.south) |- (creditoJoin) -- (creditot.north);
\draw[arrow] (creditot.east) -- (carterat.west);
\draw[arrow] (carterat.east) -- (corr.west);

\draw[sidearrow] (creditot.south) -- (creditom.north);
\draw[sidearrow] (carterat.south) -- (carteram.north);

\draw[arrow] (corr.south) |- (corrSplitL) -- (rho0t.north);
\draw[arrow] (corr.south) |- (corrSplitR) -- (rho1t.north);
\draw[sidearrow] (rho0t.south) -- (rho0m.north);
\draw[sidearrow] (rho1t.south) -- (rho1m.north);

\draw[arrow] (rho0m.south) |- (resLeftJoin) -- (reservast.west);
\draw[arrow] (rho1m.south) |- (resRightJoin) -- (reservast.east);
\draw[arrow] (reservast.south) -- (reservasm.north);
\draw[arrow] (reservasm.south) |- (0,-271) -- (fin.north);

\end{tikzpicture}
\end{adjustbox}%
}

\end{document}
